\documentclass[10pt]{article}
\usepackage[margin=20mm]{geometry}
\usepackage{amsmath,amssymb,mathtools,fontspec,unicode-math,xurl,hyperref}
\IfFontExistsTF{DejaVu Serif}{\setmainfont{DejaVu Serif}}{\setmainfont{Latin Modern Roman}}
\setmathfont{latinmodern-math.otf}
\setlength{\parindent}{0pt}\setlength{\parskip}{6pt}\sloppy
\title{A Technical Review of
The Three-Dimensional Stochastic Zakharov System}
\author{AI-assisted technical review}
\date{}
\begin{document}
\maketitle
Abstract: This review checks the conservative-noise rescaling, the ground-state coercivity argument, and the probabilistic end of the nonconservative-noise regularization argument in arXiv:2301.02089v2. The energy argument and the large-imaginary-amplitude probability limit pass the checks stated below. The unrestricted complex-amplitude reading of the limit in Theorem 1.7 is not supported by its decisive probability estimate. A fixed-imaginary-part family disproves that auxiliary limit, without disproving the final non-blow-up theorem. Because fixed-direction amplification is a plausible intended reading, the report records a parameter-scope clarification with assessment Needs Evidence. Scientific adjudication remains pending.\par
\section*{Identity, version, and coverage}
The reviewed paper is by Sebastian Herr, Michael Röckner, Martin Spitz, and Deng Zhang [paper]. The acquired file is the 46-page arXiv:2301.02089v2, submitted on 25 April 2024. The related journal identifier is 10.1214/24-AOP1710 (https://doi.org/10.1214/24-AOP1710), The Annals of Probability 53(3), 848–905 (2025). That 58-page journal text has not been compared with the acquired manuscript; the findings here apply only to the latter. PDF page numbers are one-based and agree with its printed page numbers.\par
The review reads the main statements and hypotheses on pp. 1–6, the transformations and rescaling proof on pp. 7–13, and the complete ground-state and noise-regularization arguments on pp. 31–37. The normal-form and multilinear estimates in Sections 4–5, the continuation construction in Section 6, and the appendices have not been independently re-proved. They are dependencies of the arguments checked here, not silently certified results. Source acquisition found no separate supplement or computational repository in its bounded search; the present version-metadata search likewise identified no separate correction notice. Neither search proves absence. No author code was executed.\par
\section*{Definitions and checks that support the arguments}
\subsection*{Conservative rescaling}
In the conservative problem, \(W_1=i\sum_k\phi_k\beta_k\) with real spatial coefficients and \(\mu=\frac12\sum_k\phi_k^2\). Hypothesis (H), equation (1.10), imposes square summability in \(H^3\) for the first noise, the additional directional integral condition on its gradients, and square summability in \(H^1\) for the second noise. These are not the spatially constant, complex noise assumptions used in Theorem 1.7.\par
For a regularized solution, the quadratic variation in \(d\|X\|_2^2\) contributes \(\sum_k\|\phi_kX\|_2^2\,dt\); the drift \(-\mu X\) contributes its negative. The stochastic mass term vanishes because \(\operatorname{Re}(i\phi_k|X|^2)=0\). Real-valued \(\phi_k\) is essential. This verifies the algebra of (1.2), not an independent construction of the infinite-dimensional solution.\par
The Itô product calculation gives \(u=e^{-W_1}X\) and \[e^{-W_1}\Delta(e^{W_1}u)
 =\Delta u+2\nabla W_1\cdot\nabla u+
       \bigl(\Delta W_1+\nabla W_1\cdot\nabla W_1\bigr)u .\] The last contraction is bilinear, without complex conjugation. The explicit expansion in (2.10), p. 8, resolves the possible ambiguity of the preceding \(|\nabla W_1|^2\) notation. The wave convolution in (2.6) satisfies \(dT_t=i|\nabla|T_t\,dt-i\,dW_2\), giving the stated additive-noise cancellation.\par
Equation reference: eq:aggregate. One caution in reconstructing the stochastic interchange on p. 12 is that basis indices driven by the same \(\beta_k\) are not independent. The justified aggregate estimate is \[\begin{aligned}
 \mathbb E\sup_{s\le t}\left|
  \sum_k\int_0^s\sum_j a_{jk}(r)b_j(r)\,d\beta_k(r)\right|^2
 &\le C\mathbb E\int_0^t
       \sum_k\left|\sum_j a_{jk}(r)b_j(r)\right|^2dr \notag\\
 &\le C\mathbb E\int_0^t
       \sum_k\|a_{\cdot k}(r)\|_{\ell^2}^2
                    \|b(r)\|_{\ell^2}^2dr. \end{aligned}\] With the stopped \(L^2\) bound, \(|e^{W_1}|=1\), and \(\sum_k\|\phi_k\|_\infty^2<\infty\), this recovers the finite bound \(tM\|\varphi\|_2^2\sum_k\|\phi_k\|_\infty^2\) needed there. The displayed intermediate diagonal expression on that page should not be used as a general scalar BDG inequality. This review establishes the intended aggregate bound; it does not claim a PDE counterexample to the rescaling theorem from that intermediate display.\par
\subsection*{The ground-state argument}
Write \(A=\|\nabla u\|_2^2\), \(B=\|u\|_4^4\), and \(C=\|v+|u|^2\|_2^2\). Equations (1.6) and (1.15) give \[E_Z=\frac12 A-\frac14B+\frac14C,\qquad
 K=A-\frac34B,\qquad
 E_Z+\lambda^2M-\frac13K=\frac16A+\lambda^2M+\frac14C.\] The final positive decomposition is exact. For the ground state, the usual Pohozaev identities give \(\|\nabla Q\|_2^2=3\|Q\|_2^2\) and \(\|Q\|_4^4=4\|Q\|_2^2\), hence \(4E_S(Q)M(Q)=J(Q)^2\).\par
For \(M(X_0)>0\), fixing \(\lambda_*=J(Q)/(2M(X_0))\) completes the square in the energy–mass inequality, exactly as on p. 31. The variational characterization (7.1) prevents \(K(X(t))\) from crossing zero before the energy threshold. The displayed positive decomposition then bounds \(X\) in \(H^1\) and \(Y+|X|^2\) in \(L^2\); the embedding \(H^1(\mathbb R^3)\subset L^4(\mathbb R^3)\) supplies the remaining \(Y\) bound. For zero initial Schrödinger component, \(X=0\) and the linear stochastic wave is bounded on every finite time interval. A uniform bound over an infinite time interval is unnecessary for the blow-up alternative. Thus the checked argument gives existence up to the threshold stopping time, not stochastic energy conservation or an unconditional global-existence statement.\par
\section*{Parameter scope in the regularization limit}
Problem 1 (Which noise component becomes large?).  Theorem 1.7, equation (1.17), p. 6, assumes a spatially constant complex \(\phi=\phi^{(1)}_1\) with \(\operatorname{Im}\phi\ne0\), all other first-noise coefficients zero, and \(W_2=0\). It states a probability limit as \(\phi\to\infty\). Step 3 of its proof, equations (8.44)–(8.45), p. 37, uses a limit that actually requires \(|\operatorname{Im}\phi|\to\infty\). The following check separates that exact requirement from mere growth of \(|\phi|\).\par
\subsection*{The decisive computation and all parameters}
Equation reference: eq:law. Put \(\phi=a+ib\), where \(a,b\in\mathbb R\) and \(b\ne0\). Equations (2.13)–(2.16), p. 10, yield \[\widehat\mu=\tfrac12(|\phi|^2-\phi^2)=b^2-iab,\qquad
 W_1=i\phi\beta=(ia-b)\beta,\qquad
 h(t)=e^{-2b\beta(t)-2b^2t}.\] Consequently \[
 \|h\|_{L^4(0,\infty)}^4
 =\int_0^\infty e^{-8b\beta(t)-8b^2t}\,dt
 \ \stackrel{\rm law}{=}\
 b^{-2}I,\qquad I=\int_0^\infty e^{-8B_s-8s}\,ds .\] Here \(B\) is a standard real Brownian motion; the equality is in law, not an asserted pathwise identity between different amplitudes. Brownian scaling and the change of variable \(s=b^2t\) prove [eq:law] (\#eq:law). The limit \(B_s/s\to0\) almost surely makes \(I\) finite almost surely. It is also strictly positive.\par
For the fixed threshold \(q=C_3(T,\|X_0\|_{H^1},\|Y_0\|_2)^{-4}>0\) used in (8.44), this gives \[\mathbb P(\|h\|_4^4\ge q)=\mathbb P(I\ge b^2q)\longrightarrow0
 \quad\hbox{if } |b|\longrightarrow\infty .\] The constants in (8.31), (8.41), and (8.45) depend on fixed time and initial-data norms, not on \(a\) or \(b\).\par
Now fix \(T=1\), \(X_0(x)=e^{-|x|^2}\in H^1(\mathbb R^3)\), \(Y_0=0\in L^2\), \(b=1\), and let \(\phi_n=n+i\). The data are deterministic; the noise has exactly one spatially constant mode with nonzero imaginary part; the remaining modes and the second noise vanish. These are precisely the special noise conditions of Theorem 1.7; hypothesis (H) for the conservative problem is not being imposed on this different model. Although \(|\phi_n|\to\infty\), the random variable in (8.44) has exactly the same law for every \(n\).\par
Moreover, \(\mathbb P(I\ge q)>0\) for every finite \(q\). Indeed choose \(K\) so that \(e^{8K-16}>q\). The event that \(B_1\in[-K-2,-K-1]\) and \(\sup_{0\le s\le1}|B_{1+s}-B_1|<1\) has positive probability by the nondegenerate Gaussian law and independent Brownian increments. On this event, \(B_t\le-K\) for \(1\le t\le2\), so \(I\ge\int_1^2e^{8K-16}\,dt>q\). Thus the auxiliary limit (8.44) is false along this admissible family under the unrestricted complex-modulus interpretation.\par
\subsection*{Counterevidence and the precise impact}
A natural intended convention is \(\phi=\lambda\phi_0\), \(\lambda\to+\infty\), with a fixed \(\operatorname{Im}\phi_0\ne0\). That convention does force \(|\operatorname{Im}\phi|\to\infty\) and makes the displayed argument valid. The geometric Brownian coefficient and the last line of (8.45) make this restricted interpretation particularly plausible. The report therefore does not turn the ambiguity into an unconditional assertion that Theorem 1.7 is false.\par
The issue is central to the probability step of the regularization result: smallness of \(\|h\|_4\) makes both nonlinear stopping-time conditions hold until the fixed horizon. It does not concern Theorems 1.2 or 1.5. Nor does it prove blow-up at fixed \(b\). In fact \[X^{(a+ib)}(t)
 =e^{-b^2t-bB_t}e^{ia(B_t+bt)}z^{(b)}(t),\] and the reduced equations for \(z^{(b)},v^{(b)}\) depend only on \(b\). The extra factor depending on \(a\) is spatially constant and has unit modulus. Increasing \(a\) cannot change the existence event through this transformation, but that observation alone does not determine whether its probability was already one.\par
\subsection*{Supported amendment and assessment}
An explicit condition \(|\operatorname{Im}\phi|\to\infty\), or explicit amplification of a fixed nonconservative direction, restores the verified probability limit. This is a verified amendment to Step 3, conditional on the preceding deterministic estimates and local construction; it is not an independent full proof of all those estimates. The unresolved item is the intended domain of the limit in the acquired statement and its correspondence with the journal version. The assessment is Needs Evidence, with inspection outcome scope clarification. No Major finding is claimed from an interpretation that may be narrower than the displayed notation.\par
\section*{Conclusion and state}
The checked coercivity and rescaling algebra support their stated uses, subject to the coverage limits above. The regularization proof establishes a large-nonconservative-component regime. Neither arbitrary large complex modulus nor unconditional failure of the final theorem has been established by that proof check.\par
Dimension | Recorded state\par
Issue assessment | Needs Evidence: parameter-scope clarification in Problem 1\par
Repair verification | Brownian scaling and corrected large-\(|\operatorname{Im}\phi|\) limit verified\par
Author adoption | Unknown; no author contact\par
Scientific decision | Pending; no human reviews supplied\par
Remaining scientific scope | Intended limit convention; journal-version comparison; unproved deterministic dependencies\par
\section*{Supporting checks}
The source-PDF SHA256 is `cfff510ac83ddb1ef4a87c4f579c67e701fb70007164d85660f7c4108a0f34a8`. The executable symbolic check verifies the coefficient \(\widehat\mu\), the real exponential factor, the coercive energy identity, and the fixed-\(b\) independence from \(a\). These algebraic checks are distinct from the analytic Brownian argument above. Exact code, output, source-page renders, actual viewing receipts, compiler/dependency bindings, and current file inventory are retained in the matching `verification/0491` tree.\par
The official arXiv metadata observed on 29 October 2026 reports that v2 relaxed the regularity assumption in Theorem 1.7. It does not establish a correction to the amplitude convention investigated here. All scientific reading and calculations reported here were performed in this assigned conversation. Native helpers performed only manifest hashing and public version-metadata lookup. Controller UI observations and runtime self-description are recorded separately; neither is presented as independent backend certification. Initial delivery is not independent acceptance.\par
9 S. Herr, M. Röckner, M. Spitz, and D. Zhang, The three dimensional stochastic Zakharov system, arXiv:2301.02089v2 (25 April 2024). https://arxiv.org/abs/2301.02089v2. arXiv, version history and author comments for arXiv:2301.02089, observed 29 October 2026. https://arxiv.org/abs/2301.02089.\par
\end{document}
